\subsection{Ley N° 21.663: El Único Marco Plenamente Exigible Hoy}\label{ley-n-21.663-el-uxfanico-marco-plenamente-exigible-hoy}

La \textbf{Ley N° 21.663}, Marco de Ciberseguridad e Infraestructura Crítica de la Información, fue publicada el \textbf{8 de abril de 2024} con vigencia escalonada: el grueso del cuerpo legal desde el \textbf{1 de enero de 2025}, y los Artículos 5, 8 y 9 junto al Título VII desde el \textbf{1 de marzo de 2025} . La \textbf{ANCI comenzó a funcionar el 2 de enero de 2025} .

La \textbf{Ley de Presupuestos 2026} confirma su radicación en la \textbf{Partida 32 (Ministerio de Seguridad Pública), Capítulo 06, Programa 01}, con \textbf{\$4.782.293 miles} y \textbf{dotación máxima de 40 personas} . \textbf{Cuarenta funcionarios para supervisar 1.154 entidades obligadas es un cuello de botella estructural} que explica por qué la Agencia privilegia instrucciones generales y estándares autoejecutables por sobre la fiscalización caso a caso.

El \textbf{Artículo 5} fija los criterios sustantivos copulativos para calificar como \textbf{Operador de Importancia Vital (OIV)} y el \textbf{Artículo 6} el procedimiento, con revisión \textbf{al menos cada tres años} . El primer procedimiento cerró con \textbf{1.154 OIV}: 915 en la primera etapa (Resolución Exenta N° 87, de 16 de diciembre de 2025) y \textbf{239 en la segunda, cerrada el 24 de julio de 2026} . La caída de la nómina preliminar de 1.712 a 915 ---un \textbf{46,6 \%}--- demuestra que la consulta pública tuvo efecto material.

El \textbf{D.S. N° 295/2024} del Ministerio del Interior y Seguridad Pública, publicado el \textbf{1 de marzo de 2025}, fija el esquema \textbf{«3/72/15»}: alerta temprana en \textbf{3 horas}, actualización en \textbf{72 horas ---reducida a 24 h si se afectan servicios esenciales---}, plan de acción en \textbf{7 días} e informe final en \textbf{15 días} . El \textbf{Artículo 27} define el incidente de efecto significativo por \textbf{número de afectados, duración y extensión geográfica}, criterios idénticos a los de NIS2 , y una restricción operativa determinante ordena que \textbf{los reportes excluyan datos personales}.

\subsection{Derecho Supletorio y Normativa Complementaria}\label{derecho-supletorio-y-normativa-complementaria}

Hasta el 30 de noviembre de 2026 rige la \textbf{Ley N° 19.628 de 1999}, cuya tutela es el \textbf{habeas data judicial del Artículo 16} ante el juez de letras en lo civil . La \textbf{Ley N° 21.180} de Transformación Digital del Estado es modificatoria de la Ley N° 19.880 y \textbf{su vigencia es gradual e incompleta}, con tope legal el \textbf{31 de diciembre de 2027} .

De las ocho normativas adicionales incorporadas a la comparativa (Entregable 3), dos inciden directamente sobre un proyecto de transporte: la \textbf{Ley N° 21.459} sobre delitos informáticos, que impone deberes de \textbf{conservación de datos de tráfico y de abonados}, y la \textbf{Ley N° 21.729}, publicada el \textbf{13 de febrero de 2025}, que incorpora un \textbf{Artículo 26 quáter a la Ley N° 18.168} obligando a retener registros de abonados ---IMEI, MSISDN e IMSI incluidos--- \textbf{por cinco años} . \textbf{El ordenamiento ordena acumular datos identificatorios de prácticamente toda la población mientras la Ley N° 21.719 consagra minimización y limitación del plazo de conservación}: el conflicto es concreto y verificable.

\subsection{Mapa de Organismos y Zonas Grises de Competencia}\label{mapa-de-organismos-y-zonas-grises-de-competencia}

El sistema es \textbf{policéntrico y asimétrico}: la ciberseguridad está consolidada y financiada bajo el \textbf{Ministerio de Seguridad Pública}; la protección de datos está vacante bajo el \textbf{Ministerio de Economía, Fomento y Turismo}; y el gobierno digital reside en la \textbf{Subsecretaría de Hacienda}. Chile trata \textbf{la ciberseguridad como asunto de orden público y la protección de datos como asunto de regulación económica}, y esa separación es la raíz de las zonas grises.

Se identificaron seis zonas grises adicionales (Entregable 4). La de mayor riesgo económico es \textbf{ZG-3: CMF frente a ANCI}, donde el principio de equivalencia normativa del \textbf{Artículo 37 de la Ley N° 21.663} sigue \textbf{pendiente de declaración formal}, con riesgo de doble exposición sancionatoria. \textbf{No consta convenio ni protocolo publicado de articulación entre estos órganos.}

\subsection{Mapa de organismos}\label{mapa-de-organismos}

{\def\LTcaptype{none} % do not increment counter
\begin{longtable}[]{@{}
  >{\raggedright\arraybackslash}p{(\linewidth - 2\tabcolsep) * \real{0.5000}}
  >{\raggedright\arraybackslash}p{(\linewidth - 2\tabcolsep) * \real{0.5000}}@{}}
\toprule\noalign{}
\begin{minipage}[b]{\linewidth}\raggedright
Organismo
\end{minipage} & \begin{minipage}[b]{\linewidth}\raggedright
Dependencia y función descritas en el maestro
\end{minipage} \\
\midrule\noalign{}
\endhead
\bottomrule\noalign{}
\endlastfoot
ANCI y CSIRT Nacional & Ministerio de Seguridad Pública. Seguridad, coordinación y recepción de incidentes. \\
Agencia de Protección de Datos Personales & Ministerio de Economía. Fiscalización del tratamiento de datos; el maestro la reporta no constituida al corte. \\
Secretaría de Gobierno Digital & Subsecretaría de Hacienda. Transformación digital y estándares del Estado. \\
Dirección del Trabajo & Control de jornada y límites del monitoreo laboral por GPS. \\
\end{longtable}
}

\textbf{Aplicabilidad al caso:} el maestro no acredita la calificación de Curimón como OIV. La condición de servicio esencial y la eventual designación como OIV deben distinguirse; las obligaciones adicionales de OIV no se presentan como una designación comprobada.
